La medición del cociente $e/k_B$ mediante la característica corriente-voltaje de un diodo semiconductor permite determinar experimentalmente la relación entre la carga elemental del electrón y la constante de Boltzmann, dos constantes fundamentales que intervienen en la descripción de los procesos eléctricos y térmicos. Este experimento confirma el modelo teórico del comportamiento de los semiconductores y demuestra cómo fenómenos microscópicos pueden emplearse para obtener parámetros fundamentales de la naturaleza. Asimismo, permite desarrollar habilidades en el análisis de datos experimentales, el tratamiento de incertidumbres y la comparación entre resultados experimentales y valores teóricos.